% Lösungen Einheit 5 von 5 – Vergleichen, Ordnen, Runden (zusammenbau v0.1)
\einheitenkopf{Einheit 5 von 5 · Vergleichen, Ordnen, Runden}

\erg{25}{a) Bruch}
%% TODO Lösung der Erklärzeile 26g) fehlt in der Bank
\erg{26}{a) $0{,}70$ und $0{,}42$ \quad b) $0{,}52$ \quad c) $3{,}81$ \quad d) $0{,}9$ \quad e) $12{,}53$ \quad f) $0{,}750$ \quad g) \ldots \quad h) $0{,}8$, denn $0{,}80 > 0{,}79$ \quad i) $0{,}38$; $0{,}39$; $0{,}4$; $0{,}405$ \quad j) $3{,}9$ \quad k) $(1{,}4 + 1{,}5) : 2 = 1{,}45$ \quad l) $\frac{3}{5} = 0{,}6$, also $\frac{3}{5} < 0{,}65$ \quad m) $38\,\% = 0{,}38$, also $0{,}4 > 38\,\%$ \quad n) $0{,}5^2 = 0{,}25$, also $0{,}5^2 < 0{,}5$ \quad o) $0{,}5^2 = 0{,}25$; die anderen sind $5{,}5$ und zweimal $0{,}55$}
\erg{27}{a) z. B. $0{,}75$}
\erg{28}{a) $0{,}45 < 48\,\% = 0{,}48 < 0{,}7^2 = 0{,}49 < \frac{1}{2} = 0{,}5$}
\erg{29}{a) $0{,}3 > 0{,}29$}
\erg{30}{a) Sie hat die Länge der Zahlen verglichen statt der Stellenwerte. Richtig: $0{,}30 > 0{,}28$, also ist $0{,}3$ größer. \quad b) $0{,}7$ hat $7$ Zehntel, $0{,}69$ nur $6$ Zehntel; die Zehntel entscheiden vor den Hundertsteln. \quad c) $\frac{65}{100} = \frac{13}{20}$ und $65\,\%$ \quad d) Reihenfolge: $13{,}38$ s; $13{,}4$ s; $13{,}45$ s – Ben ist am schnellsten, Can wird Dritter.}
